% ---------------------------------------------------------------------------
% preface.tex
% ---------------------------------------------------------------------------
\chapter*{Preface}
\addcontentsline{toc}{chapter}{Preface}
\markboth{Preface}{Preface}

Verification is where SystemVerilog stops being VHDL with different
punctuation. The design half of the language, covered in \refdesignbook, is a
translation exercise: a process is an \code{always} block, a generic is a
parameter. The verification half is not, because VHDL has no equivalent of
most of it. Classes, mailboxes, dynamic threads, constrained randomisation and
functional coverage do not map onto anything you have written before.

This booklet is the \textbf{Verification Cookbook}, one of four companions
covering \tool{SystemVerilog for VHDL Designers} in full. It assumes you
already know the design-side language (\refdesignbook), and it spends its
time on the part of SystemVerilog that exists purely to write better
testbenches.

\section*{What this booklet is}

The vocabulary, then a complete testbench, then the standard library built
on top of it. \Cref{ch:toolkit} introduces classes, dynamic types,
constrained randomisation, functional coverage, assertions and dynamic
threads: none of it is synthesisable, and none of it has a VHDL counterpart.
\Cref{ch:tb} assembles that vocabulary into a complete layered testbench,
from nothing, using no library at all. \Cref{ch:uvm} then shows that UVM is
the same testbench with a standard base class in front of each part.

As throughout the series, a \emph{trap} marks a mistake you are
statistically likely to make, a \emph{recipe} gives a pattern to copy, and a
plain note supplies background you can skip on a first reading.

\section*{What this booklet is not}

It is not a UVM reference manual; the Accellera UVM 1.2 user guide and the
IEEE 1800.2 LRM remain the authoritative sources for anything not covered
here. It does not discuss \tool{Verilator} or Python testbenches in any
depth: that material, including how much of this chapter's testbench survives
a change of simulator, is in \refsimulationbook. The worked CORDIC example
that this booklet's testbenches are built against is developed in full in
\refexamplebook.

\section*{Acknowledgement}

The testbenches built in this booklet are exercised, throughout, against the
CORDIC rotator of \refexamplebook: small enough that the whole testbench fits
in your head, and arithmetic enough that a golden model earns its keep.
